\documentclass[11pt]{amsart}

\usepackage[T1]{fontenc}
\usepackage{lmodern}
\usepackage{microtype}
\usepackage{amsmath,amssymb,amsthm}
\usepackage{booktabs}
\usepackage[hidelinks]{hyperref}

\hypersetup{
  pdftitle={A Counterexample to a Self-Adhesivity Conjecture of Boege, Bolt, and Studeny}
}

\newtheorem{theorem}{Theorem}
\newtheorem{lemma}[theorem]{Lemma}
\newtheorem{remark}[theorem]{Remark}

\newcommand{\sg}{\mathfrak{sg}}
\newcommand{\st}{\mathfrak{st}}
\newcommand{\M}{\mathcal{M}}
\newcommand{\A}{\mathcal{A}}

\title[Counterexample to a self-adhesivity conjecture]
{A Counterexample to a Self-Adhesivity Conjecture of
Boege, Bolt, and Studen\'y}

\date{}

\subjclass[2020]{62B10, 06A15, 05B35}
\keywords{conditional independence, semi-graphoid, structural semi-graphoid,
self-adhesivity, counterexample}

\begin{document}

\begin{abstract}
Boege, Bolt, and Studen\'y asked whether
\[
  \sg^{\diamond}(N)\cap\st(N)=\st^{\diamond}(N)
\]
holds for every finite variable set~$N$.
We give an explicit structural semi-graphoid on five variables that is
self-adhesive relative to the semi-graphoid frame but not relative to the
structural frame. This answers their Open Problem~3 in the negative and,
by the four-variable computation of Boege, Bolt, and Studen\'y, five
variables are minimal. The proof consists of an integer supermodular
witness, exact semi-graphoid closure computations, and a short linear
certificate for the failure of structural self-adhesivity.
\end{abstract}

\maketitle

\section{Statement and notation}

Let $\sg(N)$ and $\st(N)$ denote, respectively, the semi-graphoids and
structural semi-graphoids on a finite variable set~$N$, and let
$\mathfrak f^{\diamond}(N)$ denote the models self-adhesive relative to a
frame~$\mathfrak f$; we use the definitions of~\cite{BBS}.
Since $\st\subseteq\sg$, anti-extensivity and isotonicity of the
self-adhesion operator~\cite[Lemma~5.2]{BBS} give
\[
  \st^{\diamond}(N)
  \subseteq \sg^{\diamond}(N)\cap\st(N).
\]
Boege, Bolt, and Studen\'y verified equality computationally for $|N|=4$ and
asked whether it holds in general. Their question, Open Problem~3
of~\cite{BBS} and the question closing their Section~7.4, reads: ``We would
like to know whether the relation
$\sg^{\diamond}(N)\cap\st(N)=\st^{\diamond}(N)$ holds in general.''

\begin{theorem}\label{thm:main}
For $N=\{0,1,2,3,4\}$,
\[
  \st^{\diamond}(N)
  \subsetneq \sg^{\diamond}(N)\cap\st(N).
\]
Equivalently, the conjectural equality in~\cite[Open Problem~3, the question
closing Section~7.4]{BBS} fails on five variables, and this is the smallest
possible cardinality.
\end{theorem}

An elementary conditional-independence statement is written $ij\mid K$,
where $i,j\in N$ are distinct and $K\subseteq N\setminus\{i,j\}$.
For a set function $f\colon 2^N\to\mathbb R$, write
\[
  \Delta_f(ij\mid K)
  :=f(ijK)+f(K)-f(iK)-f(jK).
\]
The function $f$ is supermodular when all elementary differences are
nonnegative, and it induces the structural semi-graphoid
\[
  \M_f:=\{ij\mid K:\Delta_f(ij\mid K)=0\}.
\]
For pairwise disjoint $A,B,C$, let $[A,B\mid C]$ denote the set of all
$ij\mid K$ with $i\in A$, $j\in B$, and
$C\subseteq K\subseteq ABC\setminus\{i,j\}$.

\section{The counterexample}

Let
\begin{align*}
\M=\{&
01\mid\varnothing,
01\mid4,
01\mid23,
01\mid24,
01\mid234,
02\mid13,
03\mid4,
03\mid12,
04\mid12,\\
&12\mid\varnothing,
12\mid03,
13\mid\varnothing,
13\mid02,
14\mid\varnothing,
14\mid0,
23\mid01,\\
&24\mid\varnothing,
24\mid0,
24\mid3,
34\mid0
\}.
\end{align*}
Juxtaposition denotes union. Define the standardized integer-valued set
function $f\colon 2^N\to\mathbb Z$ by
\[
\begin{array}{c@{\qquad}l}
\toprule
f(S)&S\\
\midrule
0&|S|\leq1,\ 01,12,13,14,24\\
1&02,03,23,34\\
2&04,014,234\\
3&012,013,123,024,124,034,134\\
4&023\\
6&0123,0124\\
7&0134,0234,1234\\
12&01234\\
\bottomrule
\end{array}
\]

\begin{lemma}\label{lem:structural}
The function $f$ is supermodular and $\M_f=\M$. Hence $\M\in\st(N)$.
\end{lemma}

\begin{proof}
There are $\binom52 2^3=80$ elementary differences. Exact evaluation gives
\[
  \bigl(\#\{s:\Delta_f(s)=r\}\bigr)_{r=0}^{4}
  =(20,20,25,14,1),
\]
and the twenty zero differences are precisely the statements listed in
$\M$. Thus every elementary difference is nonnegative and $\M_f=\M$.
\end{proof}

\begin{remark}\label{rem:noncoatom}
The twenty active elementary equalities have rank $19$ in the
$26$-dimensional space of standardized set functions. Thus $f$ lies in a
seven-dimensional face of the supermodular cone. The coatoms of
$(\st(N),\subseteq)$ correspond to the extreme rays of the cone of
standardized supermodular functions, and a seven-dimensional face is not an
extreme ray. Hence this model is not a coatom, and it lies outside the
$1319$ permutational types of five-variable coatoms tested
in~\cite[Section~7.6]{BBS}.
\end{remark}

\begin{lemma}\label{lem:sg}
The model $\M$ belongs to $\sg^{\diamond}(N)$.
\end{lemma}

\begin{proof}
Fix $L\subseteq N$ and a canonical copy $b_L\colon N\to N_L$ fixing~$L$.
Put
\[
  \A_L
  :=\M\cup b_L(\M)
    \cup[N\setminus N_L,N_L\setminus N\mid L].
\]
By Remark~1 after Definition~4.3 and Lemma~4.5 of~\cite{BBS}, one
canonical copy suffices and
\[
  \sg^{\diamond}(\M\mid L)
  =\sg(\A_L)\mathbin{\downarrow}N,
\]
where $\sg(\A_L)$ is the closure under the elementary semi-graphoid rule
\[
 \{ij\mid K,\ i\ell\mid jK\}
 \quad\Longleftrightarrow\quad
 \{i\ell\mid K,\ ij\mid \ell K\}.
\]
Corollary~4.12 of~\cite{BBS} covers $|L|\leq1$, and Lemmas~4.15--4.16 there
cover the co-singletons $|L|=4$; the semi-graphoid frame satisfies the
hypotheses of these results, and of Lemma~4.5, by \cite[Lemmas~3.1, 4.10,
and~4.15]{BBS}. With the trivial case $L=N$, only $|L|\in\{2,3\}$ requires
computation. Exact saturation gives the following sizes; in every row the
$N$-marginal of the closure is exactly the twenty-statement model~$\M$.
\[
\begin{array}{c@{\quad}rr@{\qquad}c@{\quad}rr}
\toprule
L&|\A_L|&|\sg(\A_L)|&L&|\A_L|&|\sg(\A_L)|\\
\midrule
01&183&430&012&54&172\\
02&184&494&013&54&154\\
03&184&437&014&52&52\\
04&184&577&023&56&169\\
12&183&419&024&54&81\\
13&183&374&034&54&67\\
14&183&346&123&54&165\\
23&184&405&124&53&108\\
24&183&248&134&54&67\\
34&184&325&234&54&60\\
\bottomrule
\end{array}
\]
Therefore $\sg^{\diamond}(\M\mid L)=\M$ for every $L\subseteq N$, and
$\M\in\sg^{\diamond}(N)$.
\end{proof}

\begin{lemma}\label{lem:notst}
The model $\M$ does not belong to $\st^{\diamond}(N)$.
\end{lemma}

\begin{proof}
Take $L=012$, let $N'=\{0,1,2,5,6\}$, and define the copy
$b\colon N\to N'$ by
\[
  b(0)=0,\qquad b(1)=1,\qquad b(2)=2,\qquad b(3)=5,\qquad b(4)=6.
\]
Suppose that a structural adhesion exists. It is induced by a
supermodular function $g\colon2^{\{0,\ldots,6\}}\to\mathbb R$.
Write $\delta_{ij\mid K}:=\Delta_g(ij\mid K)$. Exact equality of the two
marginals and the adhesion condition imply
\[
  \delta_s=0
  \qquad
  \text{for every }s\in
  \M\cup b(\M)\cup[34,56\mid012].
\]
A coefficientwise expansion of elementary differences gives the identity
\begin{align}
&\delta_{04\mid3}
 +\delta_{01\mid235}
 +\delta_{01\mid245}
 +\delta_{05\mid234}
 +\delta_{15\mid0234}\notag\\
&\quad
 +\delta_{34\mid25}
 +\delta_{35\mid02}
 +\delta_{35\mid12}
 +\delta_{45\mid02}
 +\delta_{45\mid12}
\label{eq:certificate}\\
={}&
 \underbrace{
 \delta_{01\mid23}
 +\delta_{04\mid12}
 -\delta_{24\mid0}
 +\delta_{24\mid3}
 +\delta_{34\mid0}}_{\text{statements in }\M}
 \notag\\
&+
 \underbrace{
 \delta_{01\mid25}
 +\delta_{05\mid12}
 +\delta_{15\mid02}}_{\text{statements in }b(\M)}
 \notag\\
&+
 \underbrace{
 2\delta_{35\mid012}
 +\delta_{45\mid012}
 +\delta_{45\mid0123}}_{\text{statements in }[34,56\mid012]}.
 \notag
\end{align}
The right-hand side is zero. Every term on the left-hand side is
nonnegative by supermodularity, so all of them vanish. In particular,
$\delta_{04\mid3}=0$, and the marginal on $N$ contains $04\mid3$.
This contradicts the required marginal $\M$, because $04\mid3\notin\M$;
indeed,
\[
  \Delta_f(04\mid3)=1.
\]
\end{proof}

\begin{proof}[Proof of Theorem~\ref{thm:main}]
Lemmas~\ref{lem:structural}, \ref{lem:sg}, and~\ref{lem:notst} give
\[
  \M\in\sg^{\diamond}(N)\cap\st(N)
  \setminus\st^{\diamond}(N).
\]
For $|N|\leq3$, the two frames coincide and are fixed by self-adhesion
\cite[Corollary~4.17]{BBS}; for $|N|=4$, the asserted equality was
established in~\cite[Section~7.4]{BBS}. The case $|N|=4$ rests on the
SAT and LP computations reported there rather than on a proof, and we
inherit it as stated. Hence five variables are minimal.
\end{proof}

\section*{Exact verification}

All computations above are determined by the data printed here and are
reproducible in exact integer and rational arithmetic. From the table
defining $f$ one evaluates the $80$ elementary differences
$\Delta_f(ij\mid K)$ and compares their multiset with the histogram in
Lemma~\ref{lem:structural} and their zero set with~$\M$. The rank assertion
of Remark~\ref{rem:noncoatom} is a row reduction of the twenty
corresponding equality functionals in the $26$ coordinates indexed by the
subsets of~$N$ of size at least two. For each of the twenty sets $L$ with
$|L|\in\{2,3\}$, the set $\A_L$ is read off from~$\M$ and saturated under
the elementary semi-graphoid rule on $N\cup N_L$, which returns the sizes
tabulated in Lemma~\ref{lem:sg} and the $N$-marginal~$\M$. Finally,
identity~\eqref{eq:certificate} is an identity between integer combinations
of the $128$ values of $g$ and is checked by expanding both sides
coefficient by coefficient. No floating-point computation is involved.

\begin{thebibliography}{1}

\bibitem{BBS}
T.~Boege, J.~H. Bolt, and M.~Studen\'y,
\emph{Self-adhesivity in lattices of abstract conditional independence
models},
Discrete Appl. Math. \textbf{361} (2025), 196--225.
\href{https://doi.org/10.1016/j.dam.2024.10.006}
{doi:10.1016/j.dam.2024.10.006}.
Numbering of results and open problems above follows the authors' e-print
\href{https://arxiv.org/abs/2402.14053v3}{arXiv:2402.14053v3}.

\end{thebibliography}

\end{document}